%% Aysa X. Fan - CV (LaTeX version)
\documentclass[10.5pt,letterpaper]{article}

\usepackage[margin=0.7in]{geometry}
\usepackage{titlesec}
\usepackage{enumitem}
\usepackage{hyperref}
\usepackage{xcolor}
\usepackage{tgtermes}
\usepackage[T1]{fontenc}
\usepackage{tabularx}
\usepackage{fontawesome5}

% Colors
\definecolor{navy}{HTML}{1D3557}
\definecolor{NAVY}{HTML}{1D3557}
\definecolor{gold}{HTML}{DAA520}

% Hyperlink style
\hypersetup{
    colorlinks=true,
    linkcolor=navy,
    urlcolor=navy,
}

% Section formatting
\titleformat{\section}
  {\bfseries\color{navy}\normalsize\uppercase}
  {}{0em}{}[\vspace{-0.5em}\textcolor{navy}{\hrule height 0.75pt}\vspace{0.3em}]
\titlespacing*{\section}{0pt}{6pt}{2pt}

% Remove page numbers
\pagestyle{empty}

% Custom commands
\newcommand{\entrydate}[2]{%
  \noindent\begin{tabularx}{\textwidth}{@{}X r@{}}
    \textbf{#1} & #2
  \end{tabularx}\vspace{-2pt}%
}

\newcommand{\entrysub}[1]{\noindent\textit{#1}\par}
\newcommand{\entryplain}[1]{\noindent #1\par}

\newcommand{\pub}[5][]{%
  \noindent\hangindent=14pt #2. \textit{#3}. #4, #5.%
  \if\relax\detokenize{#1}\relax\else{ \textbf{#1}}\fi\par\vspace{2pt}%
}

% Bold "Aysa X. Fan" in author lists
\newcommand{\me}{\textbf{Aysa X. Fan}}

\begin{document}

% ─── Header ───
\begin{center}
  {\Large\bfseries\color{navy} AYSA X. FAN}\\[2pt]
  University of Illinois Urbana-Champaign\\[1pt]
  \href{mailto:xuemof2@illinois.edu}{xuemof2@illinois.edu} ~|~
  \href{https://aysa.fan}{aysa.fan} ~|~
  \href{https://scholar.google.com/citations?user=RJ9ICpYAAAAJ}{Google Scholar}\\[1pt]
  \href{https://linkedin.com/in/aysa-x-fan}{LinkedIn} ~|~
  \href{https://github.com/aysaxfan}{GitHub}
\end{center}

% ─── Research Interests ───
\section{Research Interests}
My research sits at the intersection of natural language processing (NLP), educational data mining, and computer science education. I study how NLP and large language models can support student learning, from automating the classification of learning behaviors using expert-developed coding schemes, to generating and evaluating pedagogical materials for novice programmers. I am broadly interested in how computational methods can help researchers and educators understand learning processes at scale.

% ─── Education ───
\section{Education}

\entrydate{Doctor of Philosophy, Curriculum \& Instruction, DELTA}{2019--2026}
\entrysub{University of Illinois Urbana-Champaign}
\entryplain{Advisor: Dr.\ Luc Paquette}
\entryplain{Dissertation: Comparative Study of Code-Tracing Learning Approaches}
\vspace{2pt}

\entrydate{Master of Information and Data Science}{2016--2018}
\entrysub{University of California, Berkeley}
\vspace{2pt}

\noindent\textbf{Honours Bachelor of Science}\\
\textit{University of Toronto}\\
Specialist in Mathematics and Its Applications in Finance and Economics, Major in Statistics

% ─── Skills ───
\section{Skills}

\noindent\textbf{Programming:} Python, SQL, R\\
\textbf{NLP/ML:} Large language models, prompt/context engineering, causal inference, text classification, data annotation pipeline design, crowdsourcing\\
\textbf{LLM tools:} Local LLM deployment on GPU and Apple Silicon (vLLM, SGLang, Llama.cpp). Local LLM finetuning (Unsloth, llamafactory, huggingface transformers)\\
\textbf{Dev Tools:} Docker, Git, Flask, GCP, AWS, Linux\\
\textbf{Research Methods:} Mixed-methods experimental design, educational data mining, latent profile analysis, quantitative ethnography, qualitative coding

\section{Manuscripts}

\pub{\me, Q.\ Liu, L.\ Paquette, J.\ D.\ Pinto}
    {Evaluating LLM-Based Classification of Student Debugging Strategies}
    {Journal of Educational Data Mining (Under Review)}{2025}
\pub{R.\ H.\ Zhang, \me, X.\ Lu, R.\ Zhang}
    {Resolving Cross-Dataset Annotation Conflicts through Guideline-Data Co-Reconciliation}
    {Conference of Language Modeling (Under Review)}{2026}

% ─── Publications ───
\section{Publications}



\pub[{\color{gold}\faTrophy}~Best Student Paper Nominee]
    {\me, Q.\ Liu, L.\ Paquette, J.\ D.\ Pinto}
    {Using LLM-based Filtering to Develop Reliable Coding Schemes for Rare Debugging Strategies}
    {International Conference on Quantitative Ethnography (ICQE)}{2024}

\pub{\me, R.\ A.\ Hendrawan, Y.\ Shi, Q.\ Ma}
    {Enhancing Code Tracing Question Generation with Refined Prompts in Large Language Models}
    {ACM Technical Symposium on Computer Science Education (SIGCSE)}{2024}

\pub{\me, A.\ B.\ Lekshmi Narayanan, M.\ Hassany, J.\ Ke}
    {Evaluating the Quality of Code Comments Generated by Large Language Models for Novice Programmers}
    {Leveraging Large Language Models for Next Generation Educational Technologies Workshop (LLM4Ed) at AIED}{2024}

\pub{Y.\ Zhang, J.\ D.\ Pinto, \me, L.\ Paquette}
    {Using Problem Similarity- and Order-based Weighting to Model Learner Performance in Introductory Computer Science Problems}
    {Journal of Educational Data Mining, 15(1), 63--99}{2023}

\pub{Y.\ Zhang, L.\ Paquette, J.\ D.\ Pinto, Q.\ Liu, \me}
    {Combining Latent Profile Analysis and Programming Traces to Understand Novices' Differences in Debugging}
    {Education and Information Technologies, 28(4), 4673--4701}{2023}

\pub{Y.\ Zhang, L.\ Paquette, J.\ D.\ Pinto, \me}
    {Utilizing Programming Traces to Explore and Model the Dimensions of Novices' Code-Writing Skill}
    {Computer Applications in Engineering Education, 31(4), 1041--1058}{2023}

\pub[{\color{gold}\faTrophy}~Best Student Paper Nominee]
    {J.\ D.\ Pinto, Q.\ Liu, L.\ Paquette, Y.\ Zhang, \me}
    {Investigating the Relationship Between Programming Experience and Debugging Behaviors in an Introductory Computer Science Course}
    {International Conference on Quantitative Ethnography (ICQE)}{2023}

\pub{R.\ H.\ Zhang, \me, R.\ Zhang}
    {CONENTAIL: An Entailment-based Framework for Universal Zero and Few Shot Classification with Supervised Contrastive Pretraining}
    {Conference of the European Chapter of the Association for Computational Linguistics (EACL), pp.\ 1941--1953}{2023}

\pub{\me, R.\ H.\ Zhang, L.\ Paquette, R.\ Zhang}
    {Exploring the Potential of Large Language Models in Generating Code-Tracing Questions for Introductory Programming Courses}
    {Findings of the Association for Computational Linguistics: EMNLP}{2023}

\pub{J.\ D.\ Pinto, Y.\ Zhang, L.\ Paquette, \me}
    {Investigating Elements of Student Persistence in an Introductory Computer Science Course}
    {5th Educational Data Mining in Computer Science Education (CSEDM) Workshop}{2021}

\pub{R.\ H.\ Zhang, Q.\ Liu, \me, H.\ Ji, D.\ Zeng, F.\ Cheng, D.\ Kawahara, S.\ Kurohashi}
    {Minimize Exposure Bias of Seq2Seq Models in Joint Entity and Relation Extraction}
    {Findings of the Association for Computational Linguistics: EMNLP, pp.\ 236--246}{2020}

% ─── Experience ───
\section{Experience}

\entrydate{Graduate Research Assistant}{Sep 2019 -- Present}
\entrysub{University of Illinois Urbana-Champaign, HEDS Lab (Advisor: Dr.\ Luc Paquette)}
\begin{itemize}[leftmargin=14pt, nosep, topsep=2pt]
  \item Designing a mixed-methods randomized experiment on code-tracing instruction; built and deployed a production LLM-powered feedback website (TypeScript, React, Google Cloud, Docker) on PrairieLearn for a live classroom deployment with real students
  \item Developed an LLM-based filtering pipeline to classify rare debugging strategies from interaction logs, reducing manual annotation effort by \(\sim\)75\%; built a custom annotation tool (prototyped in Python, rebuilt in TypeScript/Electron) (ICQE 2024, JEDM under review)
  \item Benchmarked seven LLMs on multi-label classification using vLLM with structured JSON output; fine-tuned Qwen3 models (0.6B--14B) using STaR with Unsloth LoRA; also fine-tuned BERT, RoBERTa, and ModernBERT as encoder baselines
  \item Investigated the use of LLMs (GPT-4-turbo, Llama2) to generate code-tracing questions and code comments for novice programmers; evaluated with BERTScore and Mann-Whitney U tests (EMNLP 2023, SIGCSE 2024, AIED 2024)
  \item Contributed to lab research analyzing longitudinal programming log data, including feature design and interpretation for studies on debugging and learner performance
\end{itemize}
\vspace{3pt}

% ─── OLD GRA VERSION (for comparison) ───
% \entrydate{Graduate Research Assistant}{Sep 2019 -- Present}
% \entrysub{University of Illinois Urbana-Champaign, HEDS Lab (Advisor: Dr.\ Luc Paquette)}
% \vspace{2pt}
% \noindent\textit{Dissertation: Comparative Study of Code-Tracing Learning Approaches (in progress)}
% \begin{itemize}[leftmargin=14pt, nosep, topsep=2pt]
%   \item Designed a three-condition mixed-methods randomized experiment (IRB-approved) comparing instructional approaches for teaching code tracing to introductory CS students
%   \item Developed a three-stage progressive curriculum (How/Why/What to Trace) with 10+ scaffolded learning activities covering tracing tables, discrepancy identification, and strategic variable selection
%   \item Built and self-hosted an LLM-powered feedback website that provides automated, context-aware tutoring responses to students during learning activities
%   \item Implemented the end-to-end learning activity on PrairieLearn, integrating pre/post-tests, the progressive curriculum, and the LLM feedback system for data collection
%   \item Designed pre-test, post-test, and survey instruments; managing participant recruitment, consent, and data collection within a single-semester timeline
% \end{itemize}
% \vspace{2pt}
% \noindent\textit{LLM-Based Classification of Debugging Strategies (ICQE 2024; JEDM under review)}
% \begin{itemize}[leftmargin=14pt, nosep, topsep=2pt]
%   \item Conducted a literature review identifying 22 candidate debugging strategies, then refined them through expert discussion into a 10-strategy codebook (4 retained for final annotation)
%   \item Led a multi-phase annotation effort (2{,}350 episodes) with iterative inter-rater reliability calibration across four phases, achieving \(\kappa > 0.75\) for all retained strategies
%   \item Designed an LLM-based filtering pipeline using GPT-4o to inflate rare strategy occurrence from \(\sim\)4\% to 27--40\% in annotation samples, reducing manual annotation effort by \(\sim\)75\%
%   \item Built a custom browser-based annotation tool displaying side-by-side code diffs with visual highlighting and error feedback for efficient expert labeling
%   \item Benchmarked seven LLMs (six open-source + Claude Sonnet 4.5) on multi-label debugging strategy classification using zero-shot prompting with structured JSON output
%   \item Developed four structurally distinct prompt strategies through iterative error analysis on 1{,}255 development episodes, incorporating hierarchical expert decision flows and contrastive definitions
%   \item Fine-tuned Qwen3 models (0.6B--14B) using Self-Taught Reasoner (STaR) with iterative self-improvement over 10 iterations; also fine-tuned BERT, RoBERTa, and ModernBERT as encoder baselines
%   \item Evaluated all models on a held-out 995-episode set using macro-averaged F1, Cohen's \(\kappa\), bootstrap confidence intervals, McNemar's tests, and Wilcoxon signed-rank tests (computed via scikit-learn)
% \end{itemize}
% \vspace{2pt}
% \noindent\textit{LLM-Generated Code-Tracing Questions (EMNLP 2023, SIGCSE 2024)}
% \begin{itemize}[leftmargin=14pt, nosep, topsep=2pt]
%   \item Curated a dataset of 176 code snippet--question pairs from CSAwesome and YouTube, pairing human-authored tracing questions with LLM-generated counterparts
%   \item Engineered and iteratively refined prompts for GPT-3.5-turbo and GPT-4, comparing zero-shot and few-shot generation; used BERTScore to assess question diversity and select the final model/approach
%   \item Designed a human evaluation framework with five criteria (relevance, clarity, difficulty, code relevance, authorship discernibility) and coordinated blind evaluations by four expert evaluators across 176 question pairs
%   \item Conducted Mann-Whitney U tests showing LLM-generated questions achieved comparable median quality to human-authored ones; built confusion matrices revealing 56\% of GPT-4 questions were mistaken as human-created
%   \item Identified quality features from literature (complexity, guessability, reverse tracing) and designed expert-guided prompts incorporating chain-of-thought self-selection, yielding improved complexity (3.5 vs.\ 2.5) and reverse-tracing coverage (83\% vs.\ 0\%)
% \end{itemize}
% \vspace{2pt}
% \noindent\textit{LLM-Generated Code Comments for Novice Programmers (LLM4Ed at AIED 2024)}
% \begin{itemize}[leftmargin=14pt, nosep, topsep=2pt]
%   \item Designed prompts for GPT-4, GPT-3.5-Turbo, and Llama2 to generate instructional code comments on 30 LeetCode Java solutions, benchmarked against expert-written comments
%   \item Developed a two-phase evaluation codebook: first-round Likert-scale criteria (8 dimensions) and refined second-round binary criteria with a ``friendly tutor'' qualitative measure; coordinated blind expert evaluations
%   \item Performed Kruskal-Wallis H tests, chi-square analyses, and Mann-Whitney U tests, finding GPT-4 achieved comparable quality to human experts with no statistically significant differences
% \end{itemize}
% \vspace{2pt}
% \noindent\textit{Lab Research on Novice Programming Behaviors}
% \begin{itemize}[leftmargin=14pt, nosep, topsep=2pt]
%   \item Contributed to lab research analyzing longitudinal programming log data, including feature design and interpretation for studies on debugging, learner persistence, and performance modeling
% \end{itemize}

\entrydate{Director of Education \& Co-founder}{Jun 2017 -- Jul 2025}
\entrysub{Novel Panda Learning Centre, Richmond Hill, ON, Canada}
\begin{itemize}[leftmargin=14pt, nosep, topsep=2pt]
  \item Co-founded a learning center offering children's art and enrichment programs; designed curricula, taught classes, and trained and supervised instructors
  \item Guided students through language, gesture, and demonstration rather than direct correction, fostering independent observation and creative self-expression
\end{itemize}
\vspace{3pt}

\entrydate{Product Quality Analyst}{May 2016 -- May 2019}
\entrysub{Kik Interactive, Waterloo, ON, Canada}
\begin{itemize}[leftmargin=14pt, nosep, topsep=2pt]
  \item Analyzed large-scale user interaction logs and clickstream data to support A/B testing and user segmentation in SQL/Python
  \item Built analytical dashboards and delivered data-driven product and competitive analyses to inform strategy
\end{itemize}
\vspace{3pt}

\entrydate{NLP Data Curation Team Lead}{Aug 2014 -- May 2016}
\entrysub{Maluuba (acquired by Microsoft, 2017; now part of Microsoft Research NLP), Waterloo, ON, Canada}
\begin{itemize}[leftmargin=14pt, nosep, topsep=2pt]
  \item Designed and iterated annotation guidelines for NLP model training, managing data pipelines across 10 languages
  \item Developed crowdsourcing tests and quality control frameworks to optimize annotation reliability for evolving NLP methods
  \item Led a team coordinating multilingual NLP data collection, cleaning, and validation for OEM projects
\end{itemize}

% ─── Awards & Honors ───
\section{Awards \& Honors}

\noindent{\color{gold}\faTrophy}~Best Student Paper Nominee, International Conference on Quantitative Ethnography (ICQE), 2024\\
{\color{gold}\faTrophy}~Best Student Paper Nominee, International Conference on Quantitative Ethnography (ICQE), 2023

\end{document}
